
\needspace{3\baselineskip}
\color{uniblue}{\subsection{Коммутационные аппараты (КА)}}
\color{black}

\begin{enumerate}

%%%%%%%%%%%%%%%%%%%%%%%%%%%%%%%%%%%%%%%%%%%%%%%%%%%%%%%%%%%%%%%%%%%%% ОБЩИЕ СВЕДЕНИЯ %%%%%%%%%%%%%%%%%%%%%%%%%%%%%%%%%%%%%%%%%%%%%%%%%%%%%%%%%%%%%%
\item Общие сведения

\begin{enumerate}

\item
Модуль КА предназначен для представления высоковольтных коммутационных аппаратов (заземляющий нож, выкатной элемент, выключатель) в составе защищаемого присоединения. В устройстве коммутационные аппараты представлены в составе функционального блока «КА».

\item
В функциональном блоке предусматриваются следующие коммутационные аппараты:
\begin{itemize}
    \item силовой выключатель;
    \item заземляющий нож;
    \item выкатной элемент.
\end{itemize}

\item
Ввод функционального блока в работу на этапе параметрирования устройства осуществляется программным переключателем SGF1 «Ввод\_функции». Введенное состояние функционального блока соответствует активному состоянию сигнала «КА: Ввод».

\item
Функциональный блок может быть выведен из работы оперативно путем активации сигнала «ОВ КА», о чем сигнализирует активное состояние сигнала «КА: Оперативный вывод».

\item
Логика управления режимами работы функционального блока выполнена в соответствии с рисунком \ref{ka:img_gen}. В таблице \ref{ka:tbl_gen} приведены параметры, необходимые для настройки функционального блока. 

\vspace{1mm}
\begin{figure}[H]
\centering
\includegraphics[width=1\textwidth,height=1\textheight,keepaspectratio]{\fbpath/130. КА ОЛ/_latex/img/1.pdf}
\captionof{figure}{Логика управления режимами КА}\label{ka:img_gen}
\end{figure}

\end{enumerate}

%%%%%%%%%%%%%%%%%%%%%%%%%%%%%%%%%%%%%%%%%%%%%%%%%%%%%%%%%%% ВЫКЛЮЧАТЕЛЬ %%%%%%%%%%%%%%%%%%%%%%%%%%%%%%%%%%%%%%%%%%%%%%%%%%%%%%%%%%
\item Выключатель (В)

\begin{enumerate}

\item
В функциональном блоке «КА» выключатель представлен функцией «В». 

\item
При приеме входных сигналов от логики отключения или функций РЗА («ЛО РЗ: Отключить», «АЧР: Отключить», «ЛО ПА: Отключить», «УРОВ: Сраб на себя»), а также и от логики оперативного управления («КП/УВ: Отключить», «Опер. отключить В» -- от внешнего КП) формируется команда «В: Отключить (реле)» на внутреннее выходное реле устройства с целью коммутации выходной цепи отключения ЭМО выключателя. При этом возможна выдача как импульсного сигнала (режим «Импульсный» программного переключателя SGF2 «Режим\_откл», время импульса определяется уставкой Т2 «Тимп\_откл»), так и длительного сигнала (режим «Длительный» программного переключателя SGF2 «Режим\_откл»). 
В режиме выдачи импульсного сигнала выдаваемый импульс (а также и выдаваемый сигнал в длительном режиме) ограничивается активизацией сигнала отключенного состояния выключателя «В: Отключен» с учетом или без учета сигналов «Работа ЭМО1» и «Работа ЭМО2» (в зависимости от положения программного переключателя SGF3 «Контр\_раб\_ЭМО»).

\item
При приеме входного сигнала от логики включения или функций РЗА («ЛВ ПА: Включить», «АПВ: Включить», «ЧАПВ: Включить»), а также и от логики оперативного управления («КП/УВ: Включить», «УВ/УВ: Включить») формируется команда «В: Включить (реле)» на внутреннее выходное реле устройства с целью коммутации выходной цепи включения ЭМВ выключателя. При этом возможна выдача как импульсного сигнала (режим «Импульсный» программного переключателя SGF4 «Режим\_вкл», время импульса определяется уставкой Т3 «Тимп\_вкл»), так и длительного сигнала (режим «Длительный» программного переключателя SGF4 «Режим\_вкл»).
В режиме выдачи импульсного сигнала выдаваемый импульс (а также и выдаваемый сигнал в длительном режиме) ограничивается активизацией сигнала включенного состояния выключателя «В: Включен» с учетом или без учета сигнала «Работа ЭМВ» (в зависимости от положения программного переключателя SGF5 «Контр\_раб\_ЭМВ»). При этом возможна задержка активации сигнала «В: Включен» (определяется уставкой Т4 «Тпрод\_вкл») для сброса импульсного или длительного сигнала на включение (для исключения ситуаций так называемого «опрокидывания» выключателя при раннем съёме команды «В: Включить (реле)» -- это характерно для некоторых видов масляных выключателей).

\item
Для снятия команд «Включить (реле)» или «Отключить (реле)» при отказе высоковольтного выключателя необходимо принудительно обесточить цепи управления и подать команду «Сброс».

\item
При выводе функции из работы (отсутствие выходного сигнала «В: Ввод») происходит сброс команд включения/отключения выключателя.

\item
В логике функции реализована блокировка на прохождение команды включения выключателя «В: Включить (реле)» при приеме сигнала аварийного отключения от защит.

\item
Логика функции формирует сигналы состояния выключателя «В: Отключен», «В: Включен», «В: Промеж. положение» и «В: Неиспр. положение», при этом сигнал «В: Неиспр. положение» выдается с выдержкой времени, задаваемой уставкой Т1 «Тнеиспр\_В».

\item
Ввод функции в работу на этапе параметрирования устройства осуществляется программным переключателем SGF1 «Ввод\_функции». Введенное состояние функции соответствует активному состоянию сигнала «В: Ввод».

\item
Алгоритм работы функции выполнен в соответствии с рисунком \ref{ka:img_v}.

\vspace{1mm}
\begin{figure}[H]
\centering
\includegraphics[width=1\textwidth,height=1\textheight,keepaspectratio]{\fbpath/130. КА ОЛ/_latex/img/v.pdf}
\captionof{figure}{Функциональная схема логики выключателя}\label{ka:img_v}
\end{figure}


\end{enumerate}
%%%%%%%%%%%%%%%%%%%%%%%%%%%%%%%%%%%%%%%%%%%%%%%%%%%%%%%%%%% КОНЕЦ ВЫКЛЮЧАТЕЛЬ %%%%%%%%%%%%%%%%%%%%%%%%%%%%%%%%%%%%%%%%%%%%%%%%%%%%%%%%%%

%%%%%%%%%%%%%%%%%%%%%%%%%%%%%%%%%%%%%%%%%%%%%%%%%%%%%%%%%%% Заземляющий нож (ЗН) %%%%%%%%%%%%%%%%%%%%%%%%%%%%%%%%%%%%%%%%%%%%%%%%%%%%%%%%%%
\item Заземляющий нож (ЗН)

\begin{enumerate}

\item
В функциональном блоке «КА» заземляющий нож представлен функцией «ЗН». 

\item
При приеме входного сигнала от логики оперативного управления («УЗН: Отключить») формируется команда «ЗН: Отключить (реле)» на внутреннее выходное реле устройства с целью коммутации выходной цепи отключения заземляющего ножа. При этом возможна выдача как импульсного сигнала (режим «Импульсный» программного переключателя SGF2 «Режим\_откл», время импульса определяется уставкой Т2 «Тимп\_откл»), так и длительного сигнала (режим «Длительный» программного переключателя SGF2 «Режим\_откл»). В режиме выдачи импульсного сигнала выдаваемый импульс (а также и выдаваемый сигнал в длительном режиме) ограничивается активизацией сигнала отключенного состояния заземляющего ножа «ЗН: Отключен».

\item
При приеме входного сигнала от логики оперативного управления («УЗН: Включить») формируется команда «ЗН: Включить (реле)» на внутреннее выходное реле устройства с целью коммутации выходной цепи включения заземляющего ножа. При этом возможна выдача как импульсного сигнала (режим «Импульсный» программного переключателя SGF4 «Режим\_вкл», время импульса определяется уставкой Т3 «Тимп\_вкл»), так и длительного сигнала (режим «Длительный» программного переключателя SGF4 «Режим\_вкл»). В режиме выдачи импульсного сигнала выдаваемый импульс (а также и выдаваемый сигнал в длительном режиме) ограничивается активизацией сигнала включенного состояния заземляющего ножа «ЗН: Включен».

\item
Для снятия команд «Включить (реле)» или «Отключить (реле)» при отказе привода заземляющего ножа необходимо принудительно обесточить цепи управления и подать команду «Сброс».

\item
При выводе функции из работы (отсутствие выходного сигнала «ЗН: Ввод») происходит сброс команд включения/отключения заземляющего ножа.

\item
Логика функции формирует сигналы состояния выключателя «ЗН: Отключен», «ЗН: Включен», «ЗН: Промеж. положение» и «ЗН: Неиспр. положение», при этом сигнал «ЗН: Неиспр. положение» выдается с выдержкой времени, задаваемой уставкой Т1 «Тнеиспр\_ЗН».

\item
Ввод функции в работу на этапе параметрирования устройства осуществляется программным переключателем SGF1 «Ввод\_функции». Введенное состояние функции соответствует активному состоянию сигнала «ЗН: Ввод». Для ввода в работу функции необходимо также наличие сигнала ввода функционального блока «КА» -- «КА: Ввод».

\item
Алгоритм работы функции выполнен в соответствии с рисунком \ref{ka:img_zn}.


\vspace{1mm}
\begin{figure}[H]
\centering
\includegraphics[width=0.95\textwidth,height=0.95\textheight,keepaspectratio]{\fbpath/130. КА ОЛ/_latex/img/zn.pdf}
\captionof{figure}{Функциональная схема логики ЗН}\label{ka:img_zn}
\end{figure}

\end{enumerate}
%%%%%%%%%%%%%%%%%%%%%%%%%%%%%%%%%%%%%%%%%%%%%%%%%%%%%%%%%%% КОНЕЦ Заземляющий нож (ЗН) %%%%%%%%%%%%%%%%%%%%%%%%%%%%%%%%%%%%%%%%%%%%%%%%%%%%%%%%%%

%%%%%%%%%%%%%%%%%%%%%%%%%%%%%%%%%%%%%%%%%%%%%%%%%%%%%%%%%%% Выкатной элемент (ВЭ) %%%%%%%%%%%%%%%%%%%%%%%%%%%%%%%%%%%%%%%%%%%%%%%%%%%%%%%%%%
\item Выкатной элемент (ВЭ)

\begin{enumerate}

\item
В функциональном блоке «КА» выкатной элемент представлен функцией «ВЭ». 

\item
При приеме входного сигнала от логики оперативного управления («УВЭ: Выкатить») формируется команда «ВЭ: Выкатить (реле)» на внутреннее выходное реле устройства с целью коммутации выходной цепи операции выкатывания ВЭ. При этом возможна выдача как импульсного сигнала (режим «Импульсный» программного переключателя SGF2 «Режим\_откл», время импульса определяется уставкой Т2 «Тимп\_откл»), так и длительного сигнала (режим «Длительный» программного переключателя SGF2 «Режим\_откл»). 
В режиме выдачи импульсного сигнала выдаваемый импульс (а также и выдаваемый сигнал в длительном режиме) ограничивается активизацией сигнала контрольного положения выкатного элемента «ВЭ: Выкачен».
\item
При приеме входного сигнала от логики оперативного управления («УВЭ: Вкатить») формируется команда «ВЭ: Вкатить (реле)» на внутреннее выходное реле устройства с целью коммутации выходной цепи операции вкатывания ВЭ. При этом возможна выдача как импульсного сигнала (режим «Импульсный» программного переключателя SGF4 «Режим\_вкл», время импульса определяется уставкой Т3 «Тимп\_вкл»), так и длительного сигнала (режим «Длительный» программного переключателя SGF4 «Режим\_вкл»).
В режиме выдачи импульсного сигнала выдаваемый импульс (а также и выдаваемый сигнал в длительном режиме) ограничивается активизацией сигнала рабочего положения выкатного элемента «ВЭ: Вкачен».
\item
Для снятия команд «Вкатить (реле)» или «Выкатить (реле)» при отказе привода выкатного элемента необходимо принудительно обесточить цепи управления и подать команду «Сброс».
\item
При выводе функции из работы (отсутствие выходного сигнала «ВЭ: Ввод») происходит сброс команд вкатить/выкатить выкатного элемента.

\item
Логика функции формирует сигналы состояния выкатного элемента «ВЭ: Выкачен», «ВЭ: Вкачен», «ВЭ: Промеж. положение» и «ВЭ: Неиспр. положение», при этом сигнал «ВЭ: Неиспр. положение» выдается с выдержкой времени, задаваемой уставкой Т1 «Тнеиспр\_ВЭ».

\item
Ввод функции в работу на этапе параметрирования устройства осуществляется программным переключателем SGF1 «Ввод\_функции». Введенное состояние функции соответствует активному состоянию сигнала «ВЭ: Ввод». Для ввода в работу функции необходимо также наличие сигнала ввода функционального блока «КА» -- «КА: Ввод».

\item
Алгоритм работы функции выполнен в соответствии с рисунком \ref{ka:img_ve}.

\vspace{1mm}
\begin{figure}[H]
\centering
\includegraphics[width=0.95\textwidth,height=0.95\textheight,keepaspectratio]{\fbpath/130. КА ОЛ/_latex/img/ve.pdf}
\captionof{figure}{Функциональная схема логики ВЭ}\label{ka:img_ve}
\end{figure}


\end{enumerate}
%%%%%%%%%%%%%%%%%%%%%%%%%%%%%%%%%%%%%%%%%%%%%%%%%%%%%%%%%%% КОНЕЦ Выкатной элемент (ВЭ) %%%%%%%%%%%%%%%%%%%%%%%%%%%%%%%%%%%%%%%%%%%%%%%%%%%%%%%%%%

\end{enumerate}